% ========================================================
% BÀI TẬP TOÁN 5 - KẾT NỐI TRI THỨC VỚI CUỘC SỐNG
% Bài 35: Ôn tập chung
% Định dạng: BTTeX Part III - Trắc nghiệm điền số ngắn
% ========================================================

% Câu 1
\begin{ex}
  Số lớn nhất có 6 chữ số khác nhau là $987\,654$. Chữ số hàng trăm nghìn là bao nhiêu?
  \shortans{9}
  \loigiai{
    Chữ số là 9.
  }
\end{ex}

% Câu 2
\begin{ex}
  Tính: $145\,000 + 55\,000$. Kết quả là bao nhiêu nghìn?
  \shortans{200}
  \loigiai{
    $145 + 55 = 200$ nghìn.
  }
\end{ex}

% Câu 3
\begin{ex}
  Rút gọn phân số $\frac{36}{48}$ về tối giản. Tổng tử số và mẫu số là bao nhiêu?
  \shortans{7}
  \loigiai{
    $\frac{36}{48} = \frac{3}{4}$. Tổng là $3 + 4 = 7$.
  }
\end{ex}

% Câu 4
\begin{ex}
  Viết số thập phân gồm $8$ đơn vị và $5$ phần trăm là $8{,}05$. Chữ số hàng phần mười là bao nhiêu?
  \shortans{0}
  \loigiai{
    Chữ số hàng phần mười là 0.
  }
\end{ex}

% Câu 5
\begin{ex}
  Tính kết quả: $23{,}5 + 16{,}5$.
  \shortans{40}
  \loigiai{
    $23{,}5 + 16{,}5 = 40$.
  }
\end{ex}

% Câu 6
\begin{ex}
  Tính kết quả: $60 - 24{,}8 - 15{,}2$.
  \shortans{20}
  \loigiai{
    $60 - 40 = 20$.
  }
\end{ex}

% Câu 7
\begin{ex}
  Tính kết quả: $2{,}5 \times 16$.
  \shortans{40}
  \loigiai{
    $2{,}5 \times 16 = 40$.
  }
\end{ex}

% Câu 8
\begin{ex}
  Tính kết quả: $48 : 1{,}2$.
  \shortans{40}
  \loigiai{
    $48 : 1{,}2 = 40$.
  }
\end{ex}

% Câu 9
\begin{ex}
  Một hình tam giác có đáy $15\text{ cm}$, chiều cao $8\text{ cm}$. Diện tích là bao nhiêu xăng-ti-mét vuông?
  \shortans{60}
  \loigiai{
    $(15 \times 8) : 2 = 60\text{ cm}^2$.
  }
\end{ex}

% Câu 10
\begin{ex}
  Một hình thang có hai đáy là $14\text{ m}$ và $10\text{ m}$, chiều cao $5\text{ m}$. Diện tích là bao nhiêu mét vuông?
  \shortans{60}
  \loigiai{
    $(14 + 10) \times 5 : 2 = 60\text{ m}^2$.
  }
\end{ex}

% Câu 11
\begin{ex}
  Một hình tròn có bán kính $10\text{ cm}$. Chu vi là bao nhiêu phần mười xăng-ti-mét? (Lấy $\pi = 3{,}14$, chu vi nhân 10).
  \shortans{628}
  \loigiai{
    Chu vi $62{,}8\text{ cm}$. Nhân 10 được 628.
  }
\end{ex}

% Câu 12
\begin{ex}
  Một hình tròn có bán kính $10\text{ cm}$. Diện tích là bao nhiêu xăng-ti-mét vuông? (Lấy $\pi = 3{,}14$).
  \shortans{314}
  \loigiai{
    $3{,}14 \times 100 = 314\text{ cm}^2$.
  }
\end{ex}

% Câu 13
\begin{ex}
  Đổi $3\text{ km}^2$ thành bao nhiêu héc-ta?
  \shortans{300}
  \loigiai{
    $3 \times 100 = 300\text{ ha}$.
  }
\end{ex}

% Câu 14
\begin{ex}
  Đổi $5\text{ ha}$ thành bao nhiêu nghìn mét vuông?
  \shortans{50}
  \loigiai{
    $50\,000\text{ m}^2 = 50$ nghìn mét vuông.
  }
\end{ex}

% Câu 15
\begin{ex}
  Làm tròn số $48{,}75$ đến hàng đơn vị ta được số nào?
  \shortans{49}
  \loigiai{
    Làm tròn lên thành 49.
  }
\end{ex}

% Câu 16
\begin{ex}
  Tìm $x$ biết: $x \times 1{,}5 = 45$.
  \shortans{30}
  \loigiai{
    $x = 45 : 1{,}5 = 30$.
  }
\end{ex}

% Câu 17
\begin{ex}
  Tìm $x$ biết: $x : 2{,}5 = 12$.
  \shortans{30}
  \loigiai{
    $x = 12 \times 2{,}5 = 30$.
  }
\end{ex}

% Câu 18
\begin{ex}
  Tính nhanh: $4{,}8 \times 25 + 4{,}8 \times 75$.
  \shortans{480}
  \loigiai{
    $4{,}8 \times 100 = 480$.
  }
\end{ex}

% Câu 19
\begin{ex}
  Một ô tô đi được $180\text{ km}$ trong $4$ giờ. Vận tốc ô tô là bao nhiêu km/h?
  \shortans{45}
  \loigiai{
    $180 : 4 = 45\text{ km/h}$.
  }
\end{ex}

% Câu 20
\begin{ex}
  Một mảnh đất hình chữ nhật có chiều dài $25\text{ m}$, chiều rộng $16\text{ m}$. Diện tích mảnh đất là bao nhiêu mét vuông?
  \shortans{400}
  \loigiai{
    $25 \times 16 = 400\text{ m}^2$.
  }
\end{ex}

% Câu 21
\begin{ex}
  Mua 5 quyển vở hết $40\,000$ đồng. Mua 12 quyển vở như thế hết bao nhiêu nghìn đồng?
  \shortans{96}
  \loigiai{
    1 quyển giá 8 nghìn. 12 quyển hết $8 \times 12 = 96$ nghìn đồng.
  }
\end{ex}

% Câu 22
\begin{ex}
  Hai số có tổng là $100$ và hiệu là $20$. Số lớn là bao nhiêu?
  \shortans{60}
  \loigiai{
    $(100 + 20) : 2 = 60$.
  }
\end{ex}

% Câu 23
\begin{ex}
  Ở câu trên, số bé là bao nhiêu?
  \shortans{40}
  \loigiai{
    $100 - 60 = 40$.
  }
\end{ex}

% Câu 24
\begin{ex}
  Tính giá trị biểu thức: $125 \times 8 - 450 : 0{,}9$.
  \shortans{500}
  \loigiai{
    $1\,000 - 500 = 500$.
  }
\end{ex}

% Câu 25
\begin{ex}
  Tìm số tự nhiên bé nhất có ba chữ số chia hết cho cả $2, 3, 5$ và $9$.
  \shortans{180}
  \loigiai{
    Chia hết cho 2 và 5 tận cùng là 0. Chia hết cho 9 tổng chữ số chia hết cho 9. Số bé nhất là 180.
  }
\end{ex}
